\subsection{Computational and graphical exercises}

\begin{enumerate}
\item For $x^2+y^2=16$, write a trigonometric parameterisation and a polar equation centred at the circle centre. Draw all three descriptions on one diagram.

\item For
\[
\frac{x^2}{25}+\frac{y^2}{9}=1,
\]
find $a$, $b$, $c$, the foci, and the eccentricity. Mark the semiaxes and both foci on a complete sketch.

\item A parabola has focus $(2,0)$ and directrix $x=-2$. Derive its Cartesian equation, find its vertex, and draw the equal-distance construction for one point.

\item For
\[
\frac{x^2}{16}-\frac{y^2}{9}=1,
\]
find the vertices, foci, eccentricity, and asymptotes. Draw both branches and label which sign of the focal-distance difference corresponds to each branch.

\item Convert $r=6\cos\theta$ to Cartesian form. Identify the centre and radius, and explain why the curve passes through the pole.

\item For
\[
r=\frac{6}{1+\frac12\cos\theta},
\]
identify $p$ and $e$, classify the conic, and compute its nearest and farthest focal distances.

\item Classify and sketch the focal conics
\[
r=\frac{3}{1+0.4\cos\theta},
\qquad
r=\frac{3}{1+\cos\theta},
\qquad
r=\frac{3}{1+1.5\cos\theta}.
\]
For the hyperbolic case, state clearly which branch is traced when $r\geq0$.

\item Sketch
\[
r=\frac{4}{1+0.6\cos(\theta-\pi/4)}.
\]
Mark the focus, periapsis direction, nearest point, and farthest point.
\end{enumerate}
